\section{UVM Testbench Implementation}

\subsection{UVM Architecture}

For the GCD module, I implemented a UVM testbench architecture. One methodology could be of implement two different agents, one for input channel and one for the output one, as they have independent handshake protocols.

However, I decided to implement a single agent, because both channels share the same clock and reset signals, and the dimension of the module itself is relatively compact.

The following diagram shows the UVM architecture that I implemented.

\begin{figure}[H]
    \centering
    \includesvg[width=\columnwidth]{diagram.svg}
    \caption{UVM testbench architecture diagram}
    \label{fig:uvm_diagram}
\end{figure}


\subsection{Architecture Overview}
The primary components are structured as follows:
\begin{itemize}
    \item \textbf{UVM Tests and Sequences:} The test layer defines the scenario (random, directed, smoke) and starts the corresponding sequences. The sequences generate randomized \texttt{gcd\_sequence\_item} transactions and pass them to the environment.
    \item \textbf{Environment (Env):} contains the Agent and the Scoreboard.
    \begin{itemize}
        \item \textbf{Agent:} manages the active stimulus and passive observation of the bus. It contains the Sequencer, Driver, and Monitor.
        \begin{itemize}
            \item \textbf{Sequencer:} takes generated sequence items from the test sequences and queues them for the driver.
            \item \textbf{Driver:} pulls sequence items from the sequencer and toggles the interface pins of the DUT according to the valid/ready handshake protocols.
            \item \textbf{Monitor:} passively observes the interface pins, reconstructing completed input and output handshakes back into transaction objects to broadcast to the scoreboard.
        \end{itemize}
        \item \textbf{Scoreboard:} contains the Reference Model and Covergroups to evaluate correctness of the DUT. 
    \end{itemize}
    \item \textbf{Interface:} External to the UVM environment, this static module connects the UVM testbench to the DUT pins and houses the SystemVerilog Assertions (SVA) and Watchdog timers.
\end{itemize}

\subsection{Stimulus Generation (Sequences and Sequence Items)}
The \texttt{gcd\_sequence\_item} forms the basic transaction payload. It contains the \texttt{WIDTH}-bit operands (\texttt{a} and \texttt{b}) and a randomizable integer, \texttt{out\_ready\_delay}. This delay variable is constrained to simulate a stalled consumer, testing the \texttt{DONE} state stability of the DUT.

\begin{minted}[frame=single,framesep=10pt,breaklines]{SystemVerilog}
class gcd_sequence_item extends uvm_sequence_item;
    `uvm_object_utils(gcd_sequence_item)

    rand bit [15:0] a;
    rand bit [15:0] b;
    rand int out_ready_delay;
    bit [15:0] gcd_out;

    // Constraint: mostly ready immediately, sometimes stalled
    constraint c_ready_delay {
        out_ready_delay dist {0 := 50, [2:5] := 50};
    }
\end{minted}

I've implemented three primary sequences:
\begin{itemize}
    \item \textbf{Smoke Sequence:} A hardcoded, simple transaction to verify basic connectivity and that the DUT can complete a handshake and output a valid result.
    \item \textbf{Directed Edge Case Sequence:} Specifically targets mathematical bounds. It injects combinations of zeros, equal operands, and maximum boundary values ($2^{WIDTH}-1$) to ensure the early exit (1-cycle latency) and edge cases operate without any protocol violations.
    \item \textbf{Constrained Random Sequence:} Generates between 250 and 500 randomized transactions to thoroughly stress the general $1+N$ subtraction algorithm.
\end{itemize}

\subsection{Agent Components}
\subsubsection{Driver}
The \texttt{gcd\_driver} retrieves sequence items from the sequencer and drives the corresponding values onto the physical pins of the interface. A critical feature of the driver is its handling of the dual handshakes. It asserts \texttt{in\_valid} and holds the input operands stable until the DUT raises \texttt{in\_ready}. Once the input handshake ends, the driver implements the \texttt{out\_ready\_delay} requested by the sequence item before asserting \texttt{out\_ready}, testing the ability of the DUT to hold \texttt{gcd\_out} stable.

\subsubsection{Monitor}
The \texttt{gcd\_monitor} passively observes the bus via a dedicated monitor clocking block. It waits for simultaneous \texttt{in\_valid} and \texttt{in\_ready} assertions to capture the incoming operands, and then waits for the \texttt{out\_valid} and \texttt{out\_ready} assertions to capture the result. The complete transaction is then written to an Analysis Port for the scoreboard to evaluate.

The following code shows the passive handshake observation of the monitor using \texttt{do...while} loops:
\newpage
\begin{minted}[frame=single,framesep=10pt,breaklines]{SystemVerilog}
// Wait for input handshake
do begin
    @(vif.mon_cb);
end while (!(vif.mon_cb.in_valid === 1'b1 && vif.mon_cb.in_ready === 1'b1));

tr.a = vif.mon_cb.a_in;
tr.b = vif.mon_cb.b_in;

// Wait for output handshake
do begin
    @(vif.mon_cb);
end while (!(vif.mon_cb.out_valid === 1'b1 && vif.mon_cb.out_ready === 1'b1));

tr.gcd_out = vif.mon_cb.gcd_out;
analysis_port.write(tr);
\end{minted}

\subsection{Scoreboard and Coverage}
The \texttt{gcd\_scoreboard} acts as both the evaluator and the coverage collector. It receives transactions from the monitor and passes the input operands to a custom Golden Reference Model. 

The Reference Model calculates the expected Greatest Common Divisor using the Modulo operator, instead of the subtractive approach of the DUT. The expected result is then compared against the \texttt{gcd\_out} of the DUT.

\begin{minted}[frame=single,framesep=10pt,breaklines]{SystemVerilog}
function int calculate_gcd(int a, int b);
    if (a == 0) return b;
    if (b == 0) return a;
    while (b != 0) begin
        int temp = b;
        b = a % b; // remainder
        a = temp;  // swap
    end
    return a;
endfunction
\end{minted}

At the same time, the scoreboard samples the \texttt{gcd\_cg} Covergroup. This covergroup ensures thorough functional coverage by tracking:
\begin{itemize}
    \item Occurrences of zero, max value, and intermediate values for both \texttt{a} and \texttt{b}.
    \item The cross-coverage between the two operands.
    \item Mathematical states, specifically whether $A > B$, $B > A$, or $A = B$.
\end{itemize}

\newpage

\begin{minted}[frame=single,framesep=10pt,breaklines]{SystemVerilog}
covergroup gcd_cg;
    option.per_instance = 1;

    cp_a: coverpoint cov_transaction.a {
        bins zero = {0};
        bins max  = {((1<<WIDTH)-1)}; // 2^(WIDTH)-1
        bins others = {[1 : ((1<<WIDTH)-2)]};
    }

    cp_b: coverpoint cov_transaction.b {
        bins zero = {0};
        bins max  = {((1<<WIDTH)-1)}; // 2^(WIDTH)-1
        bins others = {[1 : ((1<<WIDTH)-2)]};
    }

    cross_a_b: cross cp_a, cp_b;

    cp_mathematical_states: coverpoint (cov_transaction.a > cov_transaction.b) {
        bins a_greater_b = {1};
        bins b_greater_a = {0}; 
    }

    cp_equal: coverpoint (cov_transaction.a == cov_transaction.b) {
        bins a_equals_b = {1}; // GCD(A,A)
    }
endgroup
\end{minted}

\subsection{Interface, Assertions, and Watchdogs}
While the UVM environment handles the dynamic stimulus and checking, the static \texttt{gcd\_if} interface is heavily utilized to monitor protocol compliance and prevent deadlocks. This is reached through a combination of SystemVerilog Assertions (SVAs) and a dynamic procedural watchdog timer.

\subsubsection{SystemVerilog Assertions (SVAs)}
Concurrent assertions have been developed inside the uvm interface to directy monitor the \texttt{ready/valid} handshake protocols and ensure data stability during stalls. For instance, to satisfy the requirement that \texttt{gcd\_out} remains stable until the consumer is ready, the following property is asserted:

\begin{minted}[frame=single,framesep=10pt,breaklines]{SystemVerilog}
// REQ 12 output stability
property p_output_stable;
    @(posedge clk) disable iff (!rst_n || $isunknown(out_valid) || $isunknown(out_ready) || $isunknown(gcd_out))
    (out_valid && !out_ready) |=> (out_valid && $stable(gcd_out));
endproperty

assert_output_stable: assert property(p_output_stable)
    else $error("PROTOCOL VIOLATION: gcd_out changed while stalled!");
\end{minted}

To verify correct reset behavior, the following property ensures that upon a reset, the module immediately transitions to the IDLE state by raising \texttt{in\_ready} and clearing all outputs:
\begin{minted}[frame=single,framesep=10pt,breaklines]{SystemVerilog}
// REQ 6,7,8,16 reset behavior
property p_reset_behavior;
    @(posedge clk) !rst_n |=> (in_ready == 1'b1 && out_valid == 1'b0 && gcd_out == '0);
endproperty

assert_reset_behavior: assert property(p_reset_behavior) 
    else $error("RESET VIOLATION: Signals did not default to IDLE state!");
\end{minted}

To guarantee the correctness of the standard handshake, this property states that once a \texttt{valid} signal is asserted (on both the input and output channel), it cannot be de-asserted before the corresponding \texttt{ready} signal acknowledges the transfer:
\begin{minted}[frame=single,framesep=10pt,breaklines]{SystemVerilog}
// REQ 9 handshake protocol
property p_valid_no_drop(valid, ready);
    @(posedge clk) disable iff (!rst_n)
    (valid && !ready) |=> valid;
endproperty

assert_in_valid_no_drop: assert property(p_valid_no_drop(in_valid, in_ready)) 
    else $error("PROTOCOL VIOLATION: in_valid dropped before in_ready!");
assert_out_valid_no_drop: assert property(p_valid_no_drop(out_valid, out_ready))
    else $error("PROTOCOL VIOLATION: out_valid dropped before out_ready!");
\end{minted}

To confirm the DUT correctly leaves the IDLE state, this assertion verifies that \texttt{in\_ready} drops immediately following a successful input handshake:
\begin{minted}[frame=single,framesep=10pt,breaklines]{SystemVerilog}
// REQ 13 FSM start
property p_in_ready_drop;
    @(posedge clk) disable iff (!rst_n)
    (in_valid && in_ready) |=> (!in_ready);
endproperty

assert_in_ready_drop: assert property(p_in_ready_drop)
    else $error("FSM VIOLATION: in_ready did not drop after input handshake!");
\end{minted}

In a similar way, to prevent duplicate reads, this property makes sure that the module drops the \texttt{out\_valid} signal in the cycle immediately following a successful output handshake:
\begin{minted}[frame=single,framesep=10pt,breaklines]{SystemVerilog}
// REQ 15 FSM to idle
property p_out_valid_drop;
    @(posedge clk) disable iff (!rst_n)
    (out_valid && out_ready) |=> (!out_valid);
endproperty

assert_out_valid_drop: assert property(p_out_valid_drop)
    else $error("FSM VIOLATION: out_valid did not drop after output handshake!");
\end{minted}

To evaluate the protocol compliance from the producer, the next assertion monitors the input channel to strictly enforce that the operands remain perfectly stable during a stalled handshake:
\begin{minted}[frame=single,framesep=10pt,breaklines]{SystemVerilog}
// REQ 18 input stability
property p_input_stable_when_stalled;
    @(posedge clk) disable iff (!rst_n || $isunknown(in_ready) || $isunknown(in_valid) || $isunknown(a_in) || $isunknown(b_in))
    (in_valid && !in_ready) |=> (in_valid && $stable(a_in) && $stable(b_in));
endproperty

assert_input_stable: assert property(p_input_stable_when_stalled)
    else $error("PROTOCOL VIOLATION: in_valid dropped or input operands changed while stalled!");
\end{minted}

Finally, the following property can be considered as a concurrent watchdog to formally verify that the DUT always produces a valid output within the theoretical maximum latency limit:
\begin{minted}[frame=single,framesep=10pt,breaklines]{SystemVerilog}
property p_no_infinite_hang;
    @(posedge clk) disable iff (!rst_n || $isunknown(in_valid) || $isunknown(out_ready))
    (in_valid && in_ready) |=> ##[1:MAX_TIMEOUT] out_valid;
endproperty

assert_no_infinite_hang: assert property(p_no_infinite_hang)
    else $error("TIMEOUT: Module exceeded maximum cycles without asserting out_valid!");
\end{minted}

\subsubsection{Dynamic Procedural Watchdog}
To be sure that the DUT does not hang indefinitely  (during complex mathematical edge cases, important when WIDTH as a high value), I implemented a watchdog timer. At first I've used a hardcoded value, but then I developed a watchdog that dynamically calculates the maximum theoretical cycles required for convergence based on the specific inputs captured during the handshake.

If the DUT fails to assert \texttt{out\_valid} within this calculated window (plus a small margin of 10 cycles for handshake synchronization), the simulation throws a fatal error (\texttt{\$fatal}).

\begin{minted}[frame=single,framesep=10pt,breaklines]{SystemVerilog}
int count = 0;
int unsigned local_timeout = count_cyc(mon_cb.a_in, mon_cb.b_in) + 10;  

$display("[%0t] TIMEOUT LIMIT = %0d cycles", $time, local_timeout);

while (count <= local_timeout) begin      
    @(mon_cb);
    count++;
    if (mon_cb.out_valid || !rst_n)
        break;
end

if (count > local_timeout) begin          
    $fatal(1, "[TIMEOUT] DUT hung! Exceeded max theoretical cycles (%0d)", local_timeout);
end
\end{minted}

\begin{minted}[frame=single,framesep=10pt,breaklines]{SystemVerilog}
function automatic int unsigned count_cyc(bit [WIDTH-1:0] a, bit [WIDTH-1:0] b);
    int unsigned counter = 0;
    if (a == 0 || b == 0 || a == b) return 1;
    while (a != 0 && b != 0 && a != b) begin
        if (a > b) a -= b; else b -= a;
        counter++;
    end
    return counter;
endfunction
\end{minted}